% 级数（综述）
% license CCBYSA3
% type Wiki

本文根据 CC-BY-SA 协议转载翻译自维基百科\href{https://en.wikipedia.org/wiki/Series_(mathematics)}{相关文章}

在数学中，级数大致而言就是将无限多个项一个接一个地相加。[1]对级数的研究是微积分及其推广——数学分析——的重要组成部分。级数在数学的大多数领域都有应用，甚至在组合学中研究有限结构时，也会通过生成函数来使用级数。无限级数的数学性质使其在物理学、计算机科学、统计学和金融等其他定量学科中也得到了广泛应用。

在古希腊人看来，一个“潜在无限”的求和可以得到一个有限的结果，这被认为是悖论，其中最著名的例子就是芝诺悖论。[2][3]尽管如此，古希腊数学家们，包括阿基米德，仍然在实践中使用了无限级数，例如在抛物线求积问题中。[4][5]芝诺悖论中的数学问题在 17 世纪通过“极限”概念得到了化解，尤其是通过艾萨克·牛顿的早期微积分。[6]在19世纪，通过卡尔·弗里德里希·高斯和奥古斯丁-路易·柯西等人的工作，这一解决方案得到了进一步的严格化与改进。[7]他们回答了哪些和式是存在的问题（通过实数的完备性），以及级数项是否可以重新排列而不改变其和的问题（通过绝对收敛与条件收敛的概念）。

在现代术语中，任何一个有序的无限序列$(a_{1}, a_{2}, a_{3}, \ldots)$，无论这些项是数字、函数、矩阵，还是其他可以相加的对象，都定义了一个级数，即把这些$a_i$ 一个接一个地相加。为了强调项的数量是无限的，级数常常也被称为无限级数，以区别于有限级数，后者有时用来指有限和。级数通常写作：
$$
a_{1} + a_{2} + a_{3} + \cdots ,~
$$
或者使用大写希腊字母 Σ 的求和符号表示：[8]
$$
\sum_{i=1}^{\infty} a_i .~
$$
由于级数所表达的无限次加法无法在有限时间内逐一完成，但如果这些项及其有限部分和属于一个具有极限的集合，就可能为该级数赋予一个值，称为级数的和。这个值定义为：当 $n \to \infty$ 时，级数前$n$项的部分和的极限（若该极限存在）。[9][10][11] 这些有限和称为部分和。用求和符号表示：
$$
\sum_{i=1}^{\infty} a_i = \lim_{n \to \infty} \sum_{i=1}^{n} a_i ,~
$$
若极限存在。[9][10][11]当极限存在时，该级数称为收敛或可和，此时序列$(a_{1}, a_{2}, a_{3}, \ldots)$ 也是可和的；反之，如果极限不存在，则称该级数发散。[9][10][11]
